\chapter{Introduction}
\label{ch:introduction}

\section{Motivation}

Large language models trained for general-purpose assistance can be prompted into role-specific assistants, tutors, domain experts, professional agents, and produce fluent responses inside that role for a few turns at a time. They cannot be relied upon to stay inside the role across longer interactions. The model retrieves accurate information and answers the question on the table, but in doing so it contradicts a stated self-fact, flips a worldview claim, or violates a role constraint that was visible in the prompt at the start of the conversation. Wang and colleagues documented this drift empirically on LLaMA-2-70B over eight-round self-chats\autocite{wang2024drift}, and the same pattern shows up qualitatively in every long-form role-playing system we have inspected.

The standard retrieval-augmented generation pipeline, the recipe by which an external corpus is plumbed into an LLM's prompt, is not designed for this problem\autocite{lewis2020rag}. RAG conditions retrieval on the user query. It does not condition on who the assistant is supposed to be, and it does not check whether the retrieved content is compatible with the assistant's role. For a tutor persona answering a question about distributed systems, RAG returns the most semantically relevant passages from the corpus. If one of those passages happens to argue for a position the tutor has explicitly committed not to hold, RAG surfaces it anyway. The model then has to decide whether to ignore the retrieved content (defeating the point of RAG) or to incorporate it (defeating the point of the persona). Most current systems do neither cleanly.

Two lines of recent work frame the intervention space. Retrieval-side routing makes the retriever conditional on something other than the user query alone. Skill-RAG detects failure states in the LLM's hidden representations and routes retrieval through one of four learned skills\autocite{wei2026skillrag}. Probing-RAG operates similarly with self-probing\autocite{baek2025probingrag}. The claim shared across this thread is that failure states have geometric structure in the hidden representations and that structure can be exploited as a routing signal. Activation-side steering takes a different cut: extract a linear direction that represents what the persona is about and steer the residual stream toward it at generation time. Anthropic's Persona Vectors paper extracts the direction by contrastive prompting and shows that persona activations predict persona shifts before generation, enabling monitoring and real-time steering\autocite{chen2025personavectors}. The Assistant Axis line extends the same family of probes to Gemma-2-27B, Qwen 3-32B, and Llama 3.3-70B, with activation capping demonstrated as a guardrail against harmful drift in production conversations\autocite{lu2026assistantaxis}.

The architectural gap is the intersection. Skill-RAG demonstrates that internal-state signals route retrieval. Persona Vectors demonstrates that internal-state signals support inference-time intervention on persona drift. Nothing in the published literature, as of the April 2026 review that shaped this project's design, combines the two by retrieving on a persona-drift signal. The v0.3 design of Persona-RAG was an attempt to occupy that gap: a typed-memory architecture in which a persona-vector-based drift signal would gate when the system pays the cost of expensive consistency machinery.

\section{What this project found, and what it did not}
\label{sec:intro-findings}

This report tells two stories that do not match the project's original design.

The first story is a replication. Persona vectors as published\autocite{chen2025personavectors} extract on Gemma-2-9B-Instruct at 4-bit NF4 quantisation. The per-layer test \auroc{} reaches 1.000 across all three personas the project uses and across all four layers in the swept middle band. The Dubanowska shuffled-label and random-feature controls sit at chance and at 0.58 to 0.65 respectively, well below the 0.70 weak floor the protocol prescribes. The methodology travels to a backbone and a quantisation regime that prior work did not test, with the controls intact. We call this Experiment 1 throughout.

The second story is a refutation. The contrast-trained direction that Experiment 1 validates does not transfer to inference-time drift detection in multi-turn dialogue. Across two backbones (Gemma-2-9B-Instruct, Llama-3.1-8B-Instruct), two extraction regimes (prompt-scope and generation-scope), four layers per cell, three personas, and six hand-authored turns per condition, the maximum absolute projection delta between in-persona and drifting conversations is 0.103, in the wrong direction, against a pre-registered proceed threshold of 0.30. The signal we would have needed to gate on is missing by an order of magnitude. We call this Experiment 2 throughout.

Two architectural decisions follow. Mechanism M2 (persona-compatibility-filtered retrieval) loses its central operation, the persona-vector projection used as a rerank signal, and is dropped from the project. Mechanism M3 (the drift-gated hybrid) is preserved structurally but reverts to the v0.2 design where the gate is an LLM-as-judge call and the hybrid ranker simplifies from three signals to two. The headline architectural claim narrows from three mechanisms over four baselines to two mechanisms over three baselines, plus a methodological negative result.

What survives the refutation, we evaluate on a custom counterfactual-retrieval probe suite of 90 multi-turn conversations covering three role personas and three probe types, scored by MiniCheck\autocite{minicheck2024}, an adapted SYCON stance-flip metric\autocite{sycon2025}, and a three-judge PoLL panel\autocite{poll2024} of self-hosted open models with a human-validation calibration. The headline pattern is consistent across personas and probe types. The vanilla-RAG baseline B1 collapses under multi-turn pressure (PoLL persona-adherence 2.54 against the cluster's 4.5 to 4.7). The RoleGPT-level prompt baseline B2\autocite{rolegpt2023, rolellm2023}, the typed-retrieval mechanism M1, and the LLM-judge-gated mechanism M3 cluster within 0.20 points of each other on every probe type, including the Type-B counterfactual-injection probes that M1 and M3 were designed to address. Three diagnostic ablations (oracle drift gate, ID-RAG removal, 4-bit to fp16 precision sweep) confirm that the cluster is robust to gate calibration, to per-turn identity grounding, and to the precision regime the project tested. The binding constraint on persona-adherence at the 9B scale is not the retrieval architecture; it is the model's underlying capability.

\section{Contributions}

The contributions of this report are organised as follows.

\textbf{A replication of persona-vector extraction at quantised inference scale.} Experiment 1 reproduces the published persona-vectors methodology on Gemma-2-9B-Instruct at 4-bit NF4, on a previously unreported (backbone, quantisation) cell, with both the Dubanowska shuffled-label and random-feature controls in place. The replication is the project's first empirical contribution and the operational gate the rest of the design depended on.

\textbf{A robust negative result on persona-vector transfer to multi-turn drift detection.} Experiment 2 tests whether the validated direction differentiates in-persona from drifting multi-turn conversations on the same backbones and at the same scale, in both prompt-scope and generation-scope extraction regimes. The result is refutation at every (backbone, scope, layer) cell, with the maximum signal an order of magnitude below the proceed threshold and roughly half the cells showing the signal in the wrong direction. The empirical coverage (24 generation-scope cells, 12 prompt-scope cells) is the contribution. The architectural cascade (M2 dropped, M3 reverts, B3 deferred) is the response.

\textbf{An evaluation of the surviving architecture against strong baselines.} The counterfactual-retrieval probe suite, the metric stack, and the four-mechanism comparison map out where typed retrieval with per-turn identity grounding sits relative to a prompt-only baseline that follows the RoleGPT methodology. The mechanisms cluster tightly. The cluster is robust to three diagnostic ablations. The architectural headroom above the prompt baseline at the 9B scale is small and is bounded by model capability rather than by retrieval design.

\textbf{A reproducible artefact.} Every reported number traces to a \texttt{results/<run-id>/} directory in the source repository, every Hydra-resolved configuration is persisted alongside the metrics, and the pipeline reproduces from a fresh clone given the pinned dependency set in \texttt{pyproject.toml}.

\section{Report structure}

Chapter~\ref{ch:background} reviews the relevant literature on retrieval-augmented generation, internal-state-conditioned retrieval, persona-conditioning techniques, memory architectures for dialogue, and the metric and judging methodology this report adopts. Chapter~\ref{ch:architecture} describes the typed-memory schema and the persona-vector extraction pipeline that supports it. Chapter~\ref{ch:mechanisms} describes the surviving mechanism (M1) and the revised drift-gated mechanism (M3), alongside the baselines (B1, B2) the comparison runs against. Chapter~\ref{ch:methodology} describes the evaluation methodology, the counterfactual-retrieval probe suite, and the metric stack including the PoLL panel and the human-validation calibration. Chapter~\ref{ch:replication} reports Experiment 1 and Experiment 2 with their controls. Chapter~\ref{ch:mechanism-results} reports the counterfactual-probe evaluation across the four mechanisms. Chapter~\ref{ch:ablations} reports the three diagnostic ablations. Chapters~\ref{ch:discussion} through~\ref{ch:conclusion} interpret the results, document the limitations, and close.
